\scanpage{65}
La bijection inverse est celle qui associe à tt hexagone $Y$ l'ens $\tau$ des $s\in S-Y$ qui sont liés à au moins un élément de $Y$ (ils sont liés alors à exact\textsuperscript{mt} 4 éléments \inserted{opposés deux à deux} de l'hexagone, 2 sur chacun de ses triangles inscrits, et ils correspondent biunivoqu\textsuperscript{mt} aux trois diagonales de l'hexagone).

Autre procédé (qui revient au précédent, en remplaçant la triade $\tau$ par son opposée $\overline\tau$) : On prend l'ens des $s\in S-\overline\tau$ qui sont non-liés à tous les \unsure{éls|él|él.} de $\overline\tau$, \deleted{\unsure{et}} c'est \deleted{un} \inserted{le} $\unsure{=|\simeq}$ hexagone, \deleted{\unsure{est en somme associé}} que précéd\textsuperscript{mt} (rappelons que $\overline\tau$ est l'ens des él. de $S$ liés à tous les él. de $\tau$).

\uline{Ou encore} Pour une triade donnée $\unsure{\tau|I}$, et $x\in S-\tau$, regardons \deleted{\unsure{le nb}} $\operatorname{Card}\underbrace{\operatorname{Ét}(x)\cap\tau}_{\text{ens des }s\in\tau\text{ liés à }x}$ \unsure{,|]} soit $\unsure{\nu|v}_\tau(x)$.

\uline{Il y a} 6 él. tels que $\nu_\tau(x)=0$, ils forment un hexagone $h(\tau)$

9 él. tels que $\nu_\tau(x)=1$, ils forment un bitriangle $\mathrm{bit}(\tau)$

6 él. tels que $\nu_\tau(x)=2$, ils forment un hexagone $\overline h$ (dit ``opposé'' à $h$) $\overline h(\tau)$

3 él. tels que $\nu_\tau(x)=3$, ils forment une triade $\overline\tau$ (dite ``opposée'' à $\tau$).

% archive sheet number: 64

